\begin{table}[H]
\centering
\caption{Test for differential pre-trends by resource-increase group}
\label{tab:pretrends_test}
\begin{threeparttable}
\begin{tabular}{l c c}
\toprule
 & \multicolumn{2}{c}{Joint pre-trend test (\textit{p}-value)} \\
\cmidrule(lr){2-3}
 & Officers split & Patrol vehicles split \\
\midrule
\multicolumn{3}{@{}l}{\textit{Panel A. Outcomes}} \\
\quad Victimization & 0.823 & 0.048 \\
\quad Crime perception & 0.059 & 0.379 \\
\quad Investment in security measures & 0.820 & 0.157 \\
\quad Reported crime & 0.213 & 0.419 \\
\addlinespace
\multicolumn{3}{@{}l}{\textit{Panel B. Police resources}} \\
\quad Police officers & 0.616 & 0.444 \\
\quad Patrol vehicles & 0.405 & 0.157 \\
\bottomrule
\end{tabular}
\begin{tablenotes}[flushleft]
\footnotesize
\item Note: Each cell reports the \textit{p}-value of a joint \textit{F}-test that the high- and low-increase municipalities followed common pre-treatment trends. For each row we estimate, on the treated municipalities and pre-treatment observations only, an OLS event-study regression with comuna and year fixed effects, event-time indicators over event time -5 to -1 (reference -1), and their interactions with an indicator for the top-40\% resource increase (versus the bottom-40\%; the middle quintile is excluded); the reported test is the joint significance of the four interaction terms. Standard errors are clustered by comuna. The Officers and Patrol vehicles splits define the high/low groups by the increase in each resource, matching Figure~\ref{fig:figure_resources} panels (c) and (d). A \textit{p}-value above 0.05 indicates no differential pre-trend.
\end{tablenotes}
\end{threeparttable}
\end{table}
